\section{总结与延伸}

\subsection{核心要点回顾}

\begin{enumerate}
    \item \textbf{问题}：大多数人不知道自己的 AI skill 有 30\% 的时间在输出低质量结果，缺乏系统化测试手段。
    \item \textbf{方法来源}：借鉴 Karpathy 的 autoresearch，核心是``小改动 $\to$ 评估 $\to$ 保留或撤销''的自动循环。
    \item \textbf{关键设计}：用 3--6 个是/否 checklist 问题定义``什么叫好''，确保评估一致、可复现。过多会导致 overfitting。
    \item \textbf{实操流程}：6 步——下载 skill、选目标、回答 3 个问题、建立基准线、观察 dashboard、放手让 agent 跑。
    \item \textbf{实证结果}：落地页文案 skill 从 56\% 提升到 92\%，4 轮迭代，全程零手动。
    \item \textbf{普适性}：任何可量化的优化目标都适用——网站速度、cold outreach、newsletter、任何反复使用的 prompt。
\end{enumerate}

\subsection{方法论启示}

\begin{importantbox}{从靠运气到靠系统}
这篇文章最深刻的洞察不在于技术细节，而在于一个认知转变：
\begin{itemize}
    \item ``搭好就用''是第一阶段——你有了工具。
    \item ``搭好、测过、持续迭代''才是成熟阶段——你有了\textbf{系统}。
    \item Changelog 的积累让 skill 的改进成为\textbf{可传承的资产}，而非一次性的调参。当更强的模型出来，changelog 就是你交给新模型的``经验包''。
\end{itemize}
\end{importantbox}

\subsection{延伸阅读}

\begin{itemize}
    \item 原文作者 Ole Lehmann 的 X 账号：\href{https://x.com/itsolelehmann}{@itsolelehmann}
    \item 原文链接：\href{https://x.com/itsolelehmann/status/2033919415771713715}{How to 10x your Claude Skills}
    \item Andrej Karpathy 的 autoresearch 原始方法与仓库
    \item AutoResearch skill 下载：见原文末尾 Dropbox 链接或作者 GitHub
    \item 译者实践哥MinLi 的 X 账号：\href{https://x.com/MinLiBuilds}{@MinLiBuilds}
\end{itemize}

%% --- Fin del contenido --- %%

\end{document}
